\section*{第 \ref{ch:b3:computational-chemistry} 章——计算化学：哈特里–福克方法与密度泛函理论}

\begin{solution}{exo:b3:computational-chemistry:1}
$\dd E/\dd\alpha = \frac32 - \sqrt{2/\pi\alpha} = 0$ 给出 $\alpha = 8/9\pi = 0.283\,
a_0^{-2}$，而 $E = \frac32\alpha - 2\sqrt{2\alpha/\pi} = -4/3\pi = \qty{-0.4244}{\hartree}$，
比精确值 $-0.5$ 高出 15\,\%：单个高斯函数既没有尖点，也没有
$1s$ 函数的尾部。
\end{solution}

\begin{solution}{exo:b3:computational-chemistry:2}
STO-3G：每个 C 有 5 个函数（$1s$、$2s$、三个 $2p$），每个 H 有 1 个：$6 \times 5 + 6 =
36$。6-31G(d)：每个 C：1（内层）+ $2 \times 4$（劈裂价层）+ 6（$d$）$= 15$；
每个 H：2，$6 \times 15 + 6 \times 2 = 102$。
\end{solution}

\begin{solution}{exo:b3:computational-chemistry:3}
电子\omterm{def:b3:quantum-model-systems:schrodinger}{哈密顿算符}不含原子核质量：在
\omterm{def:b3:computational-chemistry:born-oppenheimer}{玻恩–奥本海默近似}下，\omterm{def:b3:computational-chemistry:pes}{势能面}因而 $r_e$ 与 $k = U''(r_e)$ 是
相同的。$\tilde\omega = \sqrt{k/\mu}/2\pi c$ 依赖于\omterm{def:b3:quantum-model-systems:oscillator}{约化质量}，而\omterm{def:b3:quantum-model-systems:oscillator}{约化质量}
加倍，比值为 $\sqrt2 = 1.414$（实测值 1.413）。
\end{solution}

\begin{solution}{exo:b3:computational-chemistry:4}
八个原子给出 $3N - 6 = 18$ 个振动：17 个为实频，一个
为虚频。海森矩阵的一个负\omterm{def:b3:quantum-model-systems:operator}{本征值}意味着一级鞍点，即过渡
态结构（此处例如对应旋转或消除）。
\end{solution}

\begin{solution}{exo:b3:computational-chemistry:5}
$\zeta = 27/16$，$E = -(27/16)^2 = \qty{-2.8477}{\hartree} = \qty{-77.49}{eV}$。
误差 $-2.8477 + 2.9034 = \qty{0.0557}{\hartree} = \qty{1.52}{eV}$，即
1.9\,\% 的总能量——但比许多反应能量都大。
\end{solution}

\begin{solution}{exo:b3:computational-chemistry:6}
$(-1 - E)(-0.5 - E) - (-0.2 - 0.3E)^2 = 0$，即 $0.91E^2 + 1.38E + 0.46 = 0$：
$E = \qty{-1.0217}{\hartree}$ 与 \qty{-0.4947}{\hartree}。是的：混入
第二个函数使能量低于 $H_{11}$，
正如\omterm{thm:b3:computational-chemistry:variational}{变分原理}所允许的。
\end{solution}

\begin{solution}{exo:b3:computational-chemistry:7}
$(102/36)^4 = 64$：约 640 分钟，约 11 小时。
\end{solution}

\begin{solution}{exo:b3:computational-chemistry:8}
$0.5782 \times 27.211 = \qty{15.73}{eV}$，比实测值高 0.30 eV。冻结
轨道：真实离子发生弛豫，能量更低，因此库普曼斯值偏高。
遗漏\omterm{def:b3:computational-chemistry:correlation}{电子相关}：中性分子（两个电子）的\omterm{def:b3:computational-chemistry:correlation}{相关
能}比离子（一个电子）更多，这使真实值偏大。
此处第一种误差占主导。
\end{solution}

\begin{solution}{exo:b3:computational-chemistry:9}
$-0.5960 - 2(-0.4666) = \qty{0.3372}{\hartree} = \qty{9.18}{eV}$，高于两个原子的能量。
RHF 波函数保留了 50\,\% 的 \ce{H+ H-}，而将一个质子与一个
氢负离子分开所需的代价，即 H 的电离能与
H 的电子亲和能之差，而能量平均已将其计入。
\end{solution}

\begin{solution}{exo:b3:computational-chemistry:10}
$E(c_1^2 + 2c_1c_2S + c_2^2) = c_1^2H_{11} + 2c_1c_2H_{12} + c_2^2H_{22}$。
对 $c_1$ 求导，并令 $\partial E/\partial c_1 = 0$：
$E(2c_1 + 2c_2S) = 2c_1H_{11} + 2c_2H_{12}$，即 $(H_{11} - E)c_1 + (H_{12} -
ES)c_2 = 0$；同理 $(H_{12} - ES)c_1 + (H_{22} - E)c_2 = 0$。
\end{solution}

\begin{solution}{exo:b3:computational-chemistry:11}
$k \approx [E(1.30) - 2E(1.35) + E(1.40)]/(0.05)^2 = 0.001417/0.0025 =
\qty{0.567}{\hartree\per\bohr\squared} = \qty{882}{N/m}$。取 $\mu = m_H/2 =
\qty{8.37e-28}{kg}$，$\omega = \sqrt{k/\mu} = \qty{1.03e15}{s^{-1}}$，并
$\tilde\omega = \omega/2\pi c = \qty{5452}{cm^{-1}}$，比实测的
$\tilde\omega_e$ 高出 24\,\%：\omterm{def:b3:computational-chemistry:basis}{最小基组} RHF 曲线过陡（缺少相关，\omterm{def:b3:computational-chemistry:basis}{基组}
过小），这正是计算频率需要标度修正的原因。
\end{solution}

\begin{solution}{exo:b3:computational-chemistry:12}
涉及氢键、仅相差几个 \unit{kJ/mol} 的能量差，需要计入相关与
色散：采用带色散校正的杂化泛函，或更好地采用
一种相关波函数方法计算最终能量，并采用极化的
三重劈裂\omterm{def:b3:computational-chemistry:basis}{基组}（包含\omterm{def:b3:computational-chemistry:split-valence}{弥散函数}）。在相同水平上优化两种构象，
以频率确认极小点（全为实频），并加入\omterm{def:b3:quantum-model-systems:oscillator}{零点
能}，用更大的计算检验\omterm{def:b3:computational-chemistry:basis}{基组}，并与一个
构象能已知的相关体系比较。
\end{solution}

\begin{solution}{pb:b3:computational-chemistry:1}
\textbf{1.} \omterm{def:b3:computational-chemistry:basis}{最小基组}中的每个\omterm{def:b3:computational-chemistry:basis}{斯莱特型轨道}都被一个固定的
由三个拟合它的高斯函数构成的收缩组合所代替。
\textbf{2.} 氦核对其电子的吸引更强，因此其 $1s$ 函数
更紧缩（指数更大）。$6.3624/3.4253 = 1.857 = (1.69/1.24)^2$。
\textbf{3.} 两个基函数、两个电子、一个占据轨道（及一个
虚轨道）。
\textbf{4.} $V_{\mathrm{nn}} = Z_{\mathrm{He}}Z_{\mathrm H}/R = 2/1.4632 =
\qty{1.3669}{\hartree}$。
\textbf{5.} 两个函数重叠很强（54\,\%）：两原子足够接近，
能够成键。
\textbf{6.} $h_{11}$ 是电子在氦函数中，同时受两个原子核作用、
不计其他电子时的能量；氦核（电荷为 2）对它的束缚
要紧得多。
\textbf{7.} $(h_{11} - \varepsilon)(h_{22} - \varepsilon) - (h_{12} - \varepsilon S)^2 = 0$：
$0.71185\varepsilon^2 + 2.79292\varepsilon + 2.44969 = 0$。
\textbf{8.} $\varepsilon = \qty{-2.600}{\hartree}$ 和 \qty{-1.324}{\hartree}。
\textbf{9.} $\mathbf h$ 忽略了两个电子之间的排斥；内层
轨道过于收缩，且能量过低。
\textbf{10.} $F_{\mu\nu} = h_{\mu\nu} + \sum_{\lambda\sigma}P_{\lambda\sigma}[(\mu\nu|\sigma\lambda) -
\frac12(\mu\lambda|\sigma\nu)]$，其中 $P_{\lambda\sigma} = 2C_{\lambda1}C_{\sigma1}$。
\textbf{11.} 每个电子受到其他电子电荷云的库仑排斥
（一般而言还包括交换作用）。
\textbf{12.} $\mathbf F$ 依赖于 $\mathbf P$，后者取决于
由 $\mathbf F$ 得到的轨道：方程是非线性的，并通过逐次
近似求解，直至自洽。
\textbf{13.} $c_1^2 + c_2^2 + 2c_1c_2S = 0.7684 + 0.0410 + 0.1906 = 1.0000$。
\textbf{14.} He：$1.5369 + 0.1906 = 1.727$；H：$0.0820 + 0.1906 = 0.273$
个电子。
\textbf{15.} 电荷：He $+0.27$，H $+0.73$。正电荷主要位于
氢原子上：\ce{HeH+} 是与氦原子结合的质子，即 \ce{He + H+}，
这与预期相符，因为氦的电离能（24.6 eV）远大于氢的（13.6 eV）。
\textbf{16.} $-\varepsilon_1 = \qty{1.633}{\hartree} = \qty{44.4}{eV}$。
\textbf{17.} $E_{\mathrm{el}} = -2.8418 - 1.3669 = \qty{-4.2087}{\hartree}$。
\textbf{18.} 三次循环达到 \qty{-2.8418}{\hartree}。每个循环给出一个
行列式，其能量是收敛的哈特里–福克
能量的上界（\omterm{thm:b3:computational-chemistry:variational}{变分原理}），迭代则不断改进它。
\textbf{19.} 具有 $\zeta = 2.0925$ 的\omterm{def:b3:computational-chemistry:basis}{基组}给出较低能量，因此它是
对该分子更好的\omterm{def:b3:computational-chemistry:basis}{基组}（\omterm{thm:b3:computational-chemistry:variational}{变分原理}）：在 \ce{HeH+} 中，氦函数
比自由原子更紧缩（自由原子的 $\zeta = 1.69$ 被
标准\omterm{def:b3:computational-chemistry:basis}{基组}沿用）。
\textbf{20.} 空的反键轨道 $\sigma^*$；其能量近似于
\ce{HeH+} 电子亲和能的负值（在如此小的\omterm{def:b3:computational-chemistry:basis}{基组}中近似较差）。
\textbf{21.} 零（无电子）；$\qty{-2.8078}{\hartree}$。
\textbf{22.} $-2.8078 - (-2.8418) = \qty{0.0340}{\hartree} = \qty{89}{kJ/mol}$。
\textbf{23.} 约为实测值的一半：\omterm{def:b3:computational-chemistry:basis}{最小基组}无法将氦的\omterm{def:b3:computational-chemistry:dft}{电子密度}
向质子方向极化（没有 $p$ 函数），因而键
过弱。（与此相比，\omterm{def:b3:quantum-model-systems:oscillator}{零点能}和 \qty{298}{K} 修正都很小。）
\textbf{24.} $-(24.5874 + 54.4178)/27.211 = \qty{-2.9034}{\hartree}$；STO-3G 原子
能量偏高 \qty{0.0956}{\hartree}（\qty{2.6}{eV}）。
\textbf{25.} 几何构型取决于能量随 $R$ 的变化方式；大部分
误差集中于内层，在各个构型处几乎相同，
因而相互抵消。
\textbf{26.} \textbf{$E(\ce{HeH+}) = \qty{-2.8418}{\hartree}$，在 RHF/STO-3G（标准
\omterm{def:b3:computational-chemistry:basis}{基组}）下，\qty{1.4632}{\bohr}}（以及经典的 \qty{-2.8607}{\hartree}，
对应 $\zeta_{\mathrm{He}} = 2.0925$）。
\end{solution}
